% Options for packages loaded elsewhere
% Options for packages loaded elsewhere
\PassOptionsToPackage{unicode}{hyperref}
\PassOptionsToPackage{hyphens}{url}
\PassOptionsToPackage{dvipsnames,svgnames,x11names}{xcolor}
%
\documentclass[
  11pt,
  letterpaper,
]{article}
\usepackage{xcolor}
\usepackage[top=0.5in,left=1in,right=1in,bottom=0.75in]{geometry}
\usepackage{amsmath,amssymb}
\setcounter{secnumdepth}{-\maxdimen} % remove section numbering
\usepackage{iftex}
\ifPDFTeX
  \usepackage[T1]{fontenc}
  \usepackage[utf8]{inputenc}
  \usepackage{textcomp} % provide euro and other symbols
\else % if luatex or xetex
  \usepackage{unicode-math} % this also loads fontspec
  \defaultfontfeatures{Scale=MatchLowercase}
  \defaultfontfeatures[\rmfamily]{Ligatures=TeX,Scale=1}
\fi
\usepackage{lmodern}
\ifPDFTeX\else
  % xetex/luatex font selection
  \setmainfont[]{Arial}
\fi
% Use upquote if available, for straight quotes in verbatim environments
\IfFileExists{upquote.sty}{\usepackage{upquote}}{}
\IfFileExists{microtype.sty}{% use microtype if available
  \usepackage[]{microtype}
  \UseMicrotypeSet[protrusion]{basicmath} % disable protrusion for tt fonts
}{}
\makeatletter
\@ifundefined{KOMAClassName}{% if non-KOMA class
  \IfFileExists{parskip.sty}{%
    \usepackage{parskip}
  }{% else
    \setlength{\parindent}{0pt}
    \setlength{\parskip}{6pt plus 2pt minus 1pt}}
}{% if KOMA class
  \KOMAoptions{parskip=half}}
\makeatother
% Make \paragraph and \subparagraph free-standing
\makeatletter
\ifx\paragraph\undefined\else
  \let\oldparagraph\paragraph
  \renewcommand{\paragraph}{
    \@ifstar
      \xxxParagraphStar
      \xxxParagraphNoStar
  }
  \newcommand{\xxxParagraphStar}[1]{\oldparagraph*{#1}\mbox{}}
  \newcommand{\xxxParagraphNoStar}[1]{\oldparagraph{#1}\mbox{}}
\fi
\ifx\subparagraph\undefined\else
  \let\oldsubparagraph\subparagraph
  \renewcommand{\subparagraph}{
    \@ifstar
      \xxxSubParagraphStar
      \xxxSubParagraphNoStar
  }
  \newcommand{\xxxSubParagraphStar}[1]{\oldsubparagraph*{#1}\mbox{}}
  \newcommand{\xxxSubParagraphNoStar}[1]{\oldsubparagraph{#1}\mbox{}}
\fi
\makeatother

\usepackage{color}
\usepackage{fancyvrb}
\newcommand{\VerbBar}{|}
\newcommand{\VERB}{\Verb[commandchars=\\\{\}]}
\DefineVerbatimEnvironment{Highlighting}{Verbatim}{commandchars=\\\{\}}
% Add ',fontsize=\small' for more characters per line
\usepackage{framed}
\definecolor{shadecolor}{RGB}{241,243,245}
\newenvironment{Shaded}{\begin{snugshade}}{\end{snugshade}}
\newcommand{\AlertTok}[1]{\textcolor[rgb]{0.68,0.00,0.00}{#1}}
\newcommand{\AnnotationTok}[1]{\textcolor[rgb]{0.37,0.37,0.37}{#1}}
\newcommand{\AttributeTok}[1]{\textcolor[rgb]{0.40,0.45,0.13}{#1}}
\newcommand{\BaseNTok}[1]{\textcolor[rgb]{0.68,0.00,0.00}{#1}}
\newcommand{\BuiltInTok}[1]{\textcolor[rgb]{0.00,0.23,0.31}{#1}}
\newcommand{\CharTok}[1]{\textcolor[rgb]{0.13,0.47,0.30}{#1}}
\newcommand{\CommentTok}[1]{\textcolor[rgb]{0.37,0.37,0.37}{#1}}
\newcommand{\CommentVarTok}[1]{\textcolor[rgb]{0.37,0.37,0.37}{\textit{#1}}}
\newcommand{\ConstantTok}[1]{\textcolor[rgb]{0.56,0.35,0.01}{#1}}
\newcommand{\ControlFlowTok}[1]{\textcolor[rgb]{0.00,0.23,0.31}{\textbf{#1}}}
\newcommand{\DataTypeTok}[1]{\textcolor[rgb]{0.68,0.00,0.00}{#1}}
\newcommand{\DecValTok}[1]{\textcolor[rgb]{0.68,0.00,0.00}{#1}}
\newcommand{\DocumentationTok}[1]{\textcolor[rgb]{0.37,0.37,0.37}{\textit{#1}}}
\newcommand{\ErrorTok}[1]{\textcolor[rgb]{0.68,0.00,0.00}{#1}}
\newcommand{\ExtensionTok}[1]{\textcolor[rgb]{0.00,0.23,0.31}{#1}}
\newcommand{\FloatTok}[1]{\textcolor[rgb]{0.68,0.00,0.00}{#1}}
\newcommand{\FunctionTok}[1]{\textcolor[rgb]{0.28,0.35,0.67}{#1}}
\newcommand{\ImportTok}[1]{\textcolor[rgb]{0.00,0.46,0.62}{#1}}
\newcommand{\InformationTok}[1]{\textcolor[rgb]{0.37,0.37,0.37}{#1}}
\newcommand{\KeywordTok}[1]{\textcolor[rgb]{0.00,0.23,0.31}{\textbf{#1}}}
\newcommand{\NormalTok}[1]{\textcolor[rgb]{0.00,0.23,0.31}{#1}}
\newcommand{\OperatorTok}[1]{\textcolor[rgb]{0.37,0.37,0.37}{#1}}
\newcommand{\OtherTok}[1]{\textcolor[rgb]{0.00,0.23,0.31}{#1}}
\newcommand{\PreprocessorTok}[1]{\textcolor[rgb]{0.68,0.00,0.00}{#1}}
\newcommand{\RegionMarkerTok}[1]{\textcolor[rgb]{0.00,0.23,0.31}{#1}}
\newcommand{\SpecialCharTok}[1]{\textcolor[rgb]{0.37,0.37,0.37}{#1}}
\newcommand{\SpecialStringTok}[1]{\textcolor[rgb]{0.13,0.47,0.30}{#1}}
\newcommand{\StringTok}[1]{\textcolor[rgb]{0.13,0.47,0.30}{#1}}
\newcommand{\VariableTok}[1]{\textcolor[rgb]{0.07,0.07,0.07}{#1}}
\newcommand{\VerbatimStringTok}[1]{\textcolor[rgb]{0.13,0.47,0.30}{#1}}
\newcommand{\WarningTok}[1]{\textcolor[rgb]{0.37,0.37,0.37}{\textit{#1}}}

\usepackage{longtable,booktabs,array}
\newcounter{none} % for unnumbered tables
\usepackage{calc} % for calculating minipage widths
% Correct order of tables after \paragraph or \subparagraph
\usepackage{etoolbox}
\makeatletter
\patchcmd\longtable{\par}{\if@noskipsec\mbox{}\fi\par}{}{}
\makeatother
% Allow footnotes in longtable head/foot
\IfFileExists{footnotehyper.sty}{\usepackage{footnotehyper}}{\usepackage{footnote}}
\makesavenoteenv{longtable}
\usepackage{graphicx}
\makeatletter
\newsavebox\pandoc@box
\newcommand*\pandocbounded[1]{% scales image to fit in text height/width
  \sbox\pandoc@box{#1}%
  \Gscale@div\@tempa{\textheight}{\dimexpr\ht\pandoc@box+\dp\pandoc@box\relax}%
  \Gscale@div\@tempb{\linewidth}{\wd\pandoc@box}%
  \ifdim\@tempb\p@<\@tempa\p@\let\@tempa\@tempb\fi% select the smaller of both
  \ifdim\@tempa\p@<\p@\scalebox{\@tempa}{\usebox\pandoc@box}%
  \else\usebox{\pandoc@box}%
  \fi%
}
% Set default figure placement to htbp
\def\fps@figure{htbp}
\makeatother





\setlength{\emergencystretch}{3em} % prevent overfull lines

\providecommand{\tightlist}{%
  \setlength{\itemsep}{0pt}\setlength{\parskip}{0pt}}



 


\makeatletter
\@ifpackageloaded{caption}{}{\usepackage{caption}}
\AtBeginDocument{%
\ifdefined\contentsname
  \renewcommand*\contentsname{Table of contents}
\else
  \newcommand\contentsname{Table of contents}
\fi
\ifdefined\listfigurename
  \renewcommand*\listfigurename{List of Figures}
\else
  \newcommand\listfigurename{List of Figures}
\fi
\ifdefined\listtablename
  \renewcommand*\listtablename{List of Tables}
\else
  \newcommand\listtablename{List of Tables}
\fi
\ifdefined\figurename
  \renewcommand*\figurename{Figure}
\else
  \newcommand\figurename{Figure}
\fi
\ifdefined\tablename
  \renewcommand*\tablename{Table}
\else
  \newcommand\tablename{Table}
\fi
}
\@ifpackageloaded{float}{}{\usepackage{float}}
\floatstyle{ruled}
\@ifundefined{c@chapter}{\newfloat{codelisting}{h}{lop}}{\newfloat{codelisting}{h}{lop}[chapter]}
\floatname{codelisting}{Listing}
\newcommand*\listoflistings{\listof{codelisting}{List of Listings}}
\makeatother
\makeatletter
\makeatother
\makeatletter
\@ifpackageloaded{caption}{}{\usepackage{caption}}
\@ifpackageloaded{subcaption}{}{\usepackage{subcaption}}
\makeatother
\usepackage{bookmark}
\IfFileExists{xurl.sty}{\usepackage{xurl}}{} % add URL line breaks if available
\urlstyle{same}
\hypersetup{
  pdftitle={Research Memo - Transportation Method Engineering},
  pdfauthor={Andrew K. Sanford},
  colorlinks=true,
  linkcolor={blue},
  filecolor={Maroon},
  citecolor={Blue},
  urlcolor={Blue},
  pdfcreator={LaTeX via pandoc}}


\title{Research Memo - Transportation Method Engineering}
\author{Andrew K. Sanford}
\date{2026-07-30}
\begin{document}
\maketitle


\textbf{Re:} {[}How is Dominant Transportation Method determined?{]}

Dominant Transportation is an engineered outcome variable that
identifies the greatest of two interactive scores (if any): Estimated
monthly rideshare usage \& estimated monthly public transit usage.

\begin{Shaded}
\begin{Highlighting}[]
\CommentTok{\# Set numeric values for transportation variables}
\NormalTok{usage\_monthly }\OtherTok{\textless{}{-}} \FunctionTok{c}\NormalTok{(}
  \StringTok{"Never"}                     \OtherTok{=} \DecValTok{0}\NormalTok{,}
  \StringTok{"0{-}1 days a month"}          \OtherTok{=} \DecValTok{1}\NormalTok{,}
  \StringTok{"2{-}4 days a month"}          \OtherTok{=} \DecValTok{3}\NormalTok{,}
  \StringTok{"4{-}8 days a month"}          \OtherTok{=} \DecValTok{6}\NormalTok{,}
  \StringTok{"8 or more days a month"}    \OtherTok{=} \DecValTok{8}
\NormalTok{  )}
\NormalTok{usage\_daily }\OtherTok{\textless{}{-}} \FunctionTok{c}\NormalTok{(}
  \StringTok{"1{-}2 rides in a typical day"}       \OtherTok{=} \FloatTok{1.5}\NormalTok{,}
  \StringTok{"3{-}4 rides in a typical day"}       \OtherTok{=} \FloatTok{3.5}\NormalTok{,}
  \StringTok{"5{-}6 rides in a typical day"}       \OtherTok{=} \FloatTok{5.5}\NormalTok{,}
  \StringTok{"7 or more rides in a typical day"} \OtherTok{=} \DecValTok{7}
\NormalTok{  )}

\CommentTok{\# Calculate highest usage transport method}
\NormalTok{survey }\OtherTok{\textless{}{-}}\NormalTok{ survey }\SpecialCharTok{|\textgreater{}}
  \FunctionTok{mutate}\NormalTok{(}
    
\CommentTok{\#   Convert 3{-}month usage variables to numeric}
    \FunctionTok{across}\NormalTok{(}\FunctionTok{c}\NormalTok{(q26\_public\_transit\_3months, q28\_rideshare\_3months), }\SpecialCharTok{\textasciitilde{}}\NormalTok{ usage\_monthly[.x]),}
    
\CommentTok{\#   Convert typical daily usage variables to numeric}
    \FunctionTok{across}\NormalTok{(}\FunctionTok{c}\NormalTok{(q27\_public\_transit\_dayuse, q29\_rideshare\_dayuse), }\SpecialCharTok{\textasciitilde{}}\NormalTok{ usage\_daily[.x]),}

\CommentTok{\#   Multiply usage over past 3 months by typical daily usage (usage count estimate)}
    \AttributeTok{public\_transit\_count =}\NormalTok{ q26\_public\_transit\_3months }\SpecialCharTok{*}\NormalTok{ q27\_public\_transit\_dayuse,}
    \AttributeTok{rideshare\_count =}\NormalTok{ q28\_rideshare\_3months }\SpecialCharTok{*}\NormalTok{ q29\_rideshare\_dayuse,}

\CommentTok{\#   Assign highest usage transport method (equal if BA)}
    \AttributeTok{highest\_transport =} \FunctionTok{case\_when}\NormalTok{(}
\NormalTok{      public\_transit\_count }\SpecialCharTok{\textgreater{}}\NormalTok{ rideshare\_count }\SpecialCharTok{\textasciitilde{}} \StringTok{"public\_transit"}\NormalTok{,}
\NormalTok{      rideshare\_count }\SpecialCharTok{\textgreater{}}\NormalTok{ public\_transit\_count }\SpecialCharTok{\textasciitilde{}} \StringTok{"rideshare"}\NormalTok{,}
\NormalTok{      public\_transit\_count }\SpecialCharTok{==} \DecValTok{0} \SpecialCharTok{\&}\NormalTok{ rideshare\_count }\SpecialCharTok{==} \DecValTok{0} \SpecialCharTok{\textasciitilde{}} \StringTok{"neither"}\NormalTok{,}
      \AttributeTok{.default =} \StringTok{"equal"}\NormalTok{ )}
\NormalTok{    )}
\end{Highlighting}
\end{Shaded}

After numericizing values from a participants' survey responses,
relevant questions were sorted into these two groups, depending on which
transportation method each question targeted. Each group then had the
corresponding values multiplied to estimate the number of monthly trips
for both transportation methods. The Dominant Transportation Method was
then assigned one of four values.

\begin{itemize}
\tightlist
\item
  If both scores have a value of 0, a participant is assigned to the
  `neither' category (indicated no rideshare or public transit use)
\item
  If both scores have a value \textgreater{} 0, but are still equal, a
  participant is assigned to the `equal' category (indicated equal usage
  of both methods)
\item
  If there was a greater value between estimated rideshare usage and
  estimated public transit usage, a participant was assigned to the
  category with higher usage.
\end{itemize}

\emph{NOTE: This code for engineering `Dominant Transportation Method'
occurs in the clean\_data.qmd script.}




\end{document}
